\section{Тест производительности}

\subsection{Методика}

Тест сравнивает реализованный алгоритм Кнута"---Морриса"---Пратта с наивным поиском.
Наивный алгоритм для каждой позиции текста последовательно проверяет совпадение
с образцом, поэтому в худшем случае работает за $O(|P| \cdot |T|)$.
КМП строит префикс-функцию и затем просматривает текст один раз, что даёт
$O(|P| + |T|)$.

Проводятся две серии экспериментов. В первой серии генерируется случайный текст над
алфавитом из 10 случайных чисел, длина образца равна 8. Во второй серии используется
худший случай для наивного поиска: образец состоит из 19 единиц и завершающей двойки,
а текст состоит только из единиц.

Время поиска замеряется с помощью \texttt{std::chrono::high\_resolution\_clock}.
Генерация данных и чтение входа в замер не входят. После каждого запуска бенчмарк
сравнивает полные списки найденных позиций для КМП и наивного алгоритма.

\subsection{Результаты}

\begin{table}[!ht]
\centering
\begin{tabular}{|r|r|r|}
\hline
\rowcolor{lightgray}
$|T|$ & КМП, мкс & Наивный поиск, мкс \\
\hline
10\,000 & 27 & 23 \\
\hline
100\,000 & 314 & 293 \\
\hline
500\,000 & 951 & 959 \\
\hline
1\,000\,000 & 1\,813 & 1\,731 \\
\hline
5\,000\,000 & 10\,324 & 9\,712 \\
\hline
\end{tabular}
\caption{Случайный текст, алфавит из 10 чисел, длина образца 8}
\end{table}

\begin{table}[!ht]
\centering
\begin{tabular}{|r|r|r|}
\hline
\rowcolor{lightgray}
$|T|$ & КМП, мкс & Наивный поиск, мкс \\
\hline
100\,000 & 1\,610 & 2\,310 \\
\hline
500\,000 & 1\,186 & 5\,644 \\
\hline
1\,000\,000 & 2\,708 & 10\,895 \\
\hline
5\,000\,000 & 11\,924 & 55\,733 \\
\hline
\end{tabular}
\caption{Худший случай для наивного поиска}
\end{table}

\begin{alltt}
# bash lab4/benchmark/run_bench.sh
=== Random text, alphabet=10, pattern_len=8 ===
n=10000:
KMP time: 27us
Naive search time: 23us
KMP matches: 0
Naive matches: 0
Results equal: yes
n=100000:
KMP time: 314us
Naive search time: 293us
KMP matches: 0
Naive matches: 0
Results equal: yes
n=500000:
KMP time: 951us
Naive search time: 959us
KMP matches: 0
Naive matches: 0
Results equal: yes
n=1000000:
KMP time: 1813us
Naive search time: 1731us
KMP matches: 0
Naive matches: 0
Results equal: yes
n=5000000:
KMP time: 10324us
Naive search time: 9712us
KMP matches: 0
Naive matches: 0
Results equal: yes

=== Worst case for naive: pattern=AAA...AB, text=AAA...A ===
n=100000:
KMP time: 1610us
Naive search time: 2310us
KMP matches: 0
Naive matches: 0
Results equal: yes
n=500000:
KMP time: 1186us
Naive search time: 5644us
KMP matches: 0
Naive matches: 0
Results equal: yes
n=1000000:
KMP time: 2708us
Naive search time: 10895us
KMP matches: 0
Naive matches: 0
Results equal: yes
n=5000000:
KMP time: 11924us
Naive search time: 55733us
KMP matches: 0
Naive matches: 0
Results equal: yes
\end{alltt}

\subsection{Анализ результатов}

На случайных данных наивный поиск часто прекращает проверку после первого или второго
числа образца, поэтому его практическое время близко к линейному. Из-за этого КМП
не получает заметного преимущества: на части запусков он немного быстрее, на части
медленнее из-за дополнительных обращений к префикс-функции.

В худшем случае разница становится существенной. При $|T| = 5\,000\,000$ наивный
алгоритм тратит \textbf{55\,733 мкс}, а КМП~--- \textbf{11\,924 мкс}, то есть
КМП быстрее примерно в \textbf{4.7 раза}. Это соответствует теоретической разнице
между $O(|P| \cdot |T|)$ и $O(|P| + |T|)$.

\pagebreak
